\documentclass[11pt,a4paper]{article}
\usepackage[margin=1in]{geometry}
\usepackage{amsmath,amssymb,amsthm}
\usepackage{algorithm}
\usepackage{algpseudocode}
\usepackage{booktabs}
\usepackage{hyperref}

\theoremstyle{definition}
\newtheorem{definition}{Definition}[section]
\newtheorem{theorem}{Theorem}[section]
\newtheorem{lemma}[theorem]{Lemma}
\newtheorem{remark}{Remark}[section]

\title{\textbf{Rigorous Mathematical Specification, Dual-Path Gradient Accumulation, and Duality of Tied Embeddings}}
\author{Team Scaibu \\ \small Email: \texttt{jaydeep.9664920749@gmail.com} \\ \small Architecture and Core Engine Team}
\date{October 2026}

\begin{document}
\maketitle

\begin{abstract}
This document presents the complete mathematical foundation and algorithmic specification of Weight Tying between input token embeddings and output unembedding heads ($W_{\text{out}} \equiv E$). In sequence modeling and autoregressive language models, input embeddings map token space to latent representations ($\mathbb{R}^V \to \mathbb{R}^D$), while output projections map latent vectors back to token logits ($\mathbb{R}^D \to \mathbb{R}^V$). Tying these matrices halves parameter storage for large vocabularies ($V \times D$) and acts as a strong inductive regularizer. We formalize the dual-path gradient accumulation theorem, formulate complete algorithmic pseudo-code matching the reference implementation, derive computational complexity, step through a numerical example, and discuss dimension scaling.
\end{abstract}

\section{Notation and Mathematical Preliminaries}

Let $H \in \mathbb{R}^{B \times D}$ denote the final Transformer layer hidden states across $B$ tokens and model dimension $D \in \mathbb{N}_{\ge 1}$.

\begin{itemize}
    \item $E \in \mathbb{R}^{V \times D}$: Shared embedding and output projection weight matrix.
    \item $V \in \mathbb{N}_{\ge 1}$: Output vocabulary size (e.g., $V = 32000$).
    \item $b \in \mathbb{R}^V$: Optional vocabulary logit bias vector.
    \item $Z \in \mathbb{R}^{B \times V}$: Output unnormalized logit matrix.
    \item $G_{\text{in}} \in \mathbb{R}^{V \times D}$: Loss gradient with respect to input embedding lookup.
    \item $G_{\text{out}} \in \mathbb{R}^{B \times V}$: Loss gradient with respect to output logits ($\frac{\partial \mathcal{L}}{\partial Z}$).
\end{itemize}

\begin{table}[h]
\centering
\caption{Summary of Mathematical Symbols for Weight Tying}
\begin{tabular}{lll}
\toprule
\textbf{Symbol} & \textbf{Type / Domain} & \textbf{Description} \\
\midrule
$H$ & $\mathbb{R}^{B \times D}$ & Final hidden states \\
$E$ & $\mathbb{R}^{V \times D}$ & Tied input/output embedding matrix \\
$b$ & $\mathbb{R}^V$ & Logit bias vector \\
$Z$ & $\mathbb{R}^{B \times V}$ & Output vocabulary logits \\
$G_{\text{total}}$ & $\mathbb{R}^{V \times D}$ & Accumulated tied gradient matrix \\
\bottomrule
\end{tabular}
\end{table}

\section{Problem Definition and Core Concepts}

\subsection{Intuition and Motivation}
In a language model with vocabulary $V = 100\text{k}$ and dimension $D = 4096$, storing separate input and output embedding matrices requires $2 \times (100\text{k} \times 4096) \times 4 \approx 3.28\text{ GB}$ of parameters. Furthermore, the semantic dot-product similarity between two word tokens in input space should match their generation similarity in output space. Tying $W_{\text{out}} = E$ enforces exact duality between lexical representation and token generation while eliminating half the embedding parameter footprint.

\subsection{Formal Mathematical Definition}
\begin{definition}[Tied Logits and Gradient Accumulation]
Given hidden states $H \in \mathbb{R}^{B \times D}$ and tied matrix $E \in \mathbb{R}^{V \times D}$:
\begin{equation}
Z = H E^T + \mathbf{1}_B b^T \label{eq:tied_forward}
\end{equation}
During backpropagation, because matrix $E$ participates in both the input embedding lookup and the output logit projection, the total gradient is the exact sum of gradients from both paths:
\begin{equation}
\frac{\partial \mathcal{L}}{\partial E} = \frac{\partial \mathcal{L}}{\partial E_{\text{in}}} + \left(\frac{\partial \mathcal{L}}{\partial Z}\right)^T H \label{eq:tied_grad}
\end{equation}
\end{definition}

\section{Algorithmic Specification}

The computational pipeline implemented in \texttt{NnAlgoWeightTying} is shown below.

\begin{algorithm}[H]
\caption{Weight Tying Forward Projection and Tied Gradient Accumulation}
\begin{algorithmic}[1]
\Require Hidden states $H \in \mathbb{R}^{B \times D}$, tied table $E \in \mathbb{R}^{V \times D}$, mode $\in \{\text{tied\_forward\_logits}, \text{tied\_gradient\_accumulation}\}$
\Require Optional bias $b \in \mathbb{R}^V$, optional input gradients $G_{\text{in}} \in \mathbb{R}^{V \times D}$, output gradients $G_{\text{out}} \in \mathbb{R}^{B \times V}$
\Ensure Output logits $Z \in \mathbb{R}^{B \times V}$ or accumulated gradient matrix $G_{\text{total}} \in \mathbb{R}^{V \times D}$
\State $B \gets \text{rows}(H), \quad D \gets \text{cols}(H), \quad V \gets \text{rows}(E)$
\If{$\text{mode} = \text{tied\_forward\_logits}$}
    \State $Z \gets H \cdot E^T \in \mathbb{R}^{B \times V}$
    \If{$b \text{ is not null}$}
        \For{$i \gets 1 \textbf{ to } B$}
            \For{$j \gets 1 \textbf{ to } V$}
                \State $Z[i][j] \gets Z[i][j] + b[j]$
            \EndFor
        \EndFor
    \EndIf
    \State \Return $\{\text{logits}: Z\}$
\Else
    \State $G_{\text{total}} \gets \text{new Matrix of shape } (V, D)$
    \For{$v \gets 1 \textbf{ to } V$}
        \For{$d \gets 1 \textbf{ to } D$}
            \State $\text{sum} \gets 0.0$
            \For{$b \gets 1 \textbf{ to } B$}
                \State $\text{sum} \gets \text{sum} + G_{\text{out}}[b][v] \cdot H[b][d]$
            \EndFor
            \State $G_{\text{total}}[v][d] \gets \text{sum} + (G_{\text{in}}[v][d] \text{ if } G_{\text{in}} \text{ not null else } 0.0)$
        \EndFor
    \EndFor
    \State \Return $\{\text{tied\_gradients}: G_{\text{total}}\}$
\EndIf
\end{algorithmic}
\end{algorithm}

\section{Theoretical Guarantees and Parameter Reduction}

\begin{theorem}[Exact Parameter Halving and Dual Gradient Flow]
Weight tying reduces total lexical parameter complexity from $2 V D$ to $V D$ parameters (a 50\% reduction). Furthermore, by the chain rule of calculus:
\begin{equation}
\nabla_E \mathcal{L} = \nabla_{E_{\text{in}}} \mathcal{L} + G_{\text{out}}^T H
\end{equation}
guaranteeing exact gradient accumulation under reverse-mode automatic differentiation.
\end{theorem}

\subsection{Limitations}
\begin{itemize}
    \item \textbf{Softmax Logit Variance}: Since input embeddings are typically normalized, multiplying unnormalized hidden states by $E^T$ can produce large logits, requiring scaling by $1/\sqrt{D}$ or RMSNorm before the head.
\end{itemize}

\section{Complexity Analysis}

\subsection{Time and Space Complexity}
\begin{itemize}
    \item \textbf{Time}: Forward logits computation takes $\Theta(B \cdot V \cdot D)$ FLOPs. Backward gradient contraction takes $\Theta(B \cdot V \cdot D)$ FLOPs.
    \item \textbf{Space}: $\Theta(B \cdot V)$ memory for output logit tensor.
\end{itemize}

\section{Worked Numerical Example}

Consider $B=1, D=2, V=2$, $H = [1.0, 2.0]$, $E = \begin{bmatrix} 0.5 & -0.5 \\ 1.0 & 1.0 \end{bmatrix}$, $b = [0.1, -0.1]$.

\begin{enumerate}
    \item $Z_1 = 1.0(0.5) + 2.0(-0.5) + 0.1 = 0.5 - 1.0 + 0.1 = -0.40$.
    \item $Z_2 = 1.0(1.0) + 2.0(1.0) - 0.1 = 1.0 + 2.0 - 0.1 = 2.90$.
    \item Output Logits: $Z = [-0.40, 2.90]$.
\end{enumerate}

\section{Numerical Considerations and Edge Cases}

\begin{itemize}
    \item \textbf{Optimizer State Deduplication}: In distributed training (FSDP / Megatron-LM), tied parameters must be registered as a single parameter pointer to prevent duplicate optimizer memory allocation ($m, v$).
\end{itemize}

\begin{thebibliography}{9}
\bibitem{press2016} O.~Press and L.~Wolf, ``Using the output embedding to improve language models,'' in \emph{EACL}, pp.~157--163, 2017.
\bibitem{inan2016} H.~Inan, K.~Khosravi, and R.~Socher, ``Tying word vectors and word classifiers: A loss framework for language modeling,'' in \emph{ICLR}, 2017.
\end{thebibliography}

\end{document}
